% ===========================================================================
%  六维推进映射
% ===========================================================================

\section{正交摆座的六维力与力矩映射}
\label{sec:mapping}

\subsection{推进模型系与几何}

本节全部公式写在推进模型系 $\{P\}$ 中。$x'$ 沿名义推力轴，前推进器位于
$\bm r_f=[a,0,0]^T$，尾推进器位于 $\bm r_t=[-b,0,0]^T$。
前摆绕 $z'$ 轴，尾摆绕 $y'$ 轴，因而单位推力方向分别为
\begin{equation}
\bm n_f=
\begin{bmatrix}\cos\delta_f&\sin\delta_f&0\end{bmatrix}^{T},
\qquad
\bm n_t=
\begin{bmatrix}\cos\delta_t&0&-\sin\delta_t\end{bmatrix}^{T}.
\label{eq:gimbal_rotations}
\end{equation}
这里的“前/尾”是几何位置，“上摆/下摆”是实机悬停构型名称；滚转、俯仰和偏航
只在指定飞控坐标系后使用。

两台推进器对机体的力为 $\bm F_f=T_f\bm n_f$ 和
$\bm F_t=T_t\bm n_t$。令前桨为 CW、尾桨为 CCW，并以 $Q_f,Q_t\geq0$ 表示
反扭矩幅值，则反扭矩矢量分别为 $-Q_f\bm n_f$ 与 $+Q_t\bm n_t$。
每台推进器的总力矩满足
\begin{equation}
\bm M_i=\bm r_i\times\bm F_i+\bm Q_i.
\label{eq:propulsor_moment}
\end{equation}

\subsection{完整六维映射}

将两台推进器的贡献相加，得到
\begin{equation}
\boxed{
\begin{aligned}
F_{x'} &=T_f\cos\delta_f+T_t\cos\delta_t,\\
F_{y'} &=T_f\sin\delta_f,\\
F_{z'} &=-T_t\sin\delta_t,\\
M_{x'} &=-Q_f\cos\delta_f+Q_t\cos\delta_t,\\
M_{y'} &=-bT_t\sin\delta_t-Q_f\sin\delta_f,\\
M_{z'} &= aT_f\sin\delta_f-Q_t\sin\delta_t.
\end{aligned}}
\label{eq:six_component}
\end{equation}
式\eqref{eq:six_component}与固件 \texttt{TandemVec\_Propulsion.h::computeWrench}
逐项对应。数据驱动实现还用 Rodrigues 公式从电机几何表生成推力方向，并以中心差分构造
Jacobian；宿主机回归要求其与解析矩阵逐元素相符。

\subsection{主效能与交叉项}

在 $\delta_f=\delta_t=0$、$\omega_f=\omega_t=\omega_0$ 附近，
三个主效能为
\begin{equation}
\Delta\omega\rightarrow M_{x'},\qquad
\delta_t\rightarrow M_{y'},\qquad
\delta_f\rightarrow M_{z'}.
\label{eq:model_primary_channels}
\end{equation}
但式\eqref{eq:six_component}同时给出两个一阶反扭矩投影：
$-Q_f\delta_f$ 进入 $M_{y'}$，$-Q_t\delta_t$ 进入 $M_{z'}$。
差速在非零摆角处也通过推力和反扭矩幅值改变其余力矩。因此，“正交摆轴”只保证
推力力臂主项的几何分离，不保证有限摆角处的完整 Jacobian 对角化。

\subsection{符号检查与飞控语义}

当 $\delta_f>0$ 时，$F_{y'}>0$ 且推力力臂主项 $M_{z'}>0$；
当 $\delta_t>0$ 时，$F_{z'}<0$ 且 $M_{y'}<0$。零摆角、等转速时
$M_{x'}=M_{y'}=M_{z'}=0$。这些检查只验证推进模型系的物理符号。
固件中的滚转/俯仰/偏航通道还需通过\cref{eq:frame_permutation,eq:firmware_axis_command}
转换并叠加标定符号。实机确认后的控制语义为上摆滚转、下摆俯仰、差速偏航，
不再用“前摆偏航、差速滚转”作为固件描述。

\begin{figure}[H]
\centering
\resizebox{0.98\linewidth}{!}{%
\begin{tikzpicture}[>=Stealth,
  body/.style={draw=black!65, fill=black!6, rounded corners=2pt, line width=.7pt},
  axis/.style={->, draw=black!70, line width=.6pt, >={Stealth[length=2.4mm]}},
  front/.style={->, draw=blue!70!black, line width=1.2pt, >={Stealth[length=3mm]}},
  tail/.style={->, draw=orange!85!black, line width=1.2pt, >={Stealth[length=3mm]}},
  baseline/.style={->, densely dashed, draw=black!40, line width=.65pt, >={Stealth[length=2mm]}},
  moment/.style={->, draw=teal!70!black, line width=1pt, >={Stealth[length=2.5mm]}},
  torque/.style={->, densely dashed, draw=orange!85!black, line width=.9pt, >={Stealth[length=2.2mm]}},
  spin/.style={->, draw=purple!70!black, line width=.9pt, >={Stealth[length=2.2mm]}},
  dim/.style={<->, draw=black!45, line width=.55pt, >={Stealth[length=1.7mm]}},
  annot/.style={font=\footnotesize\sffamily, fill=white, inner sep=1.2pt},
]
  % 仰视图：只有前摆在 y 方向产生分量
  \begin{scope}[shift={(0,0)}]
    \draw[body] (-2.7,-1.25) rectangle (2.7,1.25);
    \draw[axis] (-3.0,0) -- (3.25,0) node[annot, right=1pt] {$x_b$};
    \draw[axis] (0,-1.55) -- (0,1.8) node[annot, above=1pt] {$y_b$};
    \fill (0,0) circle (2pt); \node[annot, anchor=north west] at (.08,-.06) {CG};
    \node[annot] at (-2.25,.85) {$\odot\,z_b$};

    \fill[orange!85!black] (-1.65,0) circle (2.8pt);
    \fill[blue!70!black] (1.85,0) circle (2.8pt);
    \node[annot] at (-1.65,-.45) {尾电机};
    \node[annot] at (1.85,-.45) {前电机};
    \draw[spin] (-1.65,-1.42) -- (-2.45,-1.42)
      node[annot, below=1pt, text=purple!70!black] {$\bm h_t\parallel -x_b$};
    \draw[spin] (1.85,-1.42) -- (2.65,-1.42)
      node[annot, below=1pt, text=purple!70!black] {$\bm h_f\parallel +x_b$};

    \draw[baseline] (1.85,0) -- ++(0:1.65cm);
    \draw[front] (1.85,0) -- ++(32:1.85cm) node[annot, text=blue!70!black, above left=-1pt] {$\bm{T}_f$};
    \draw[blue!70!black, line width=.8pt] ($(1.85,0)+(0:7mm)$) arc (0:32:7mm);
    \node[annot, text=blue!70!black] at (2.72,.30) {$\delta_f$};

    \draw[tail] (-1.65,0) -- ++(0:1.55cm)
      node[annot, text=orange!85!black, midway, above=3pt] {$\bm{T}_t$};
    \draw[moment] (-.55,.86) arc (150:400:4.6mm);
    \node[annot, text=teal!70!black] at (-.7,1.48) {$M_z^{(T)}=aT_f\sin\delta_f$};

    \draw[dim] (-1.65,-1.05) -- node[annot] {$b$} (0,-1.05);
    \draw[dim] (0,-1.05) -- node[annot] {$a$} (1.85,-1.05);
    \node[font=\footnotesize\sffamily] at (0,-2.05) {(a) 仰视（沿$-z_b$）：前摆偏航主项};
  \end{scope}

  % 侧视图：尾部机体受力向 +x、-z
  \begin{scope}[shift={(7.4,0)}]
    \draw[body] (-2.7,-1.25) rectangle (2.7,1.25);
    \draw[axis] (-3.0,0) -- (3.25,0) node[annot, right=1pt] {$x_b$};
    \draw[axis] (0,1.8) -- (0,-1.65) node[annot, below=1pt] {$z_b$};
    \fill (0,0) circle (2pt); \node[annot, anchor=north west] at (.08,-.06) {CG};
    \node[annot, align=center] at (2.25,.92) {$\odot\,y_b$\\[-1pt]\scriptsize 视向 $-y_b$};

    \fill[orange!85!black] (-1.65,0) circle (2.8pt);
    \fill[blue!70!black] (1.85,0) circle (2.8pt);
    \node[annot] at (-1.65,-.45) {尾电机};
    \node[annot] at (1.85,-.45) {前电机};
    \draw[spin] (-1.65,-1.42) -- (-2.45,-1.42)
      node[annot, below=1pt, text=purple!70!black] {$\bm h_t\parallel -x_b$};
    \draw[spin] (1.85,-1.42) -- (2.65,-1.42)
      node[annot, below=1pt, text=purple!70!black] {$\bm h_f\parallel +x_b$};

    \draw[baseline] (-1.65,0) -- ++(0:1.65cm);
    \draw[tail] (-1.65,0) -- ++(25:1.85cm)
      node[annot, text=orange!85!black, midway, above left=-1pt] {$\bm{T}_t$};
    \draw[orange!85!black, line width=.8pt] ($(-1.65,0)+(0:7mm)$) arc (0:25:7mm);
    \node[annot, text=orange!85!black] at (-.76,.28) {$\delta_t$};
    \draw[front] (1.85,0) -- ++(0:1.55cm) node[annot, text=blue!70!black, above=2pt] {$\bm{T}_f$};

    \draw[moment] (.62,.82) arc (35:-235:4.2mm and 3.2mm);
    \node[annot, text=teal!70!black] at (.2,1.48) {$M_y^{(T)}=-bT_t\sin\delta_t$};
    \draw[dim] (-1.65,-1.05) -- node[annot] {$b$} (0,-1.05);
    \draw[dim] (0,-1.05) -- node[annot] {$a$} (1.85,-1.05);
    \node[font=\footnotesize\sffamily] at (0,-2.05) {(b) 侧视（沿$-y_b$）：尾摆俯仰主项};
  \end{scope}
  \node[annot, text=orange!85!black] at (3.7,-2.55)
    {滚转通道：$M_x=-\tau_f\cos\delta_f+\tau_t\cos\delta_t$};
\end{tikzpicture}%
}
\caption{按Web core与Three.js摆座语义重绘的推力矢量映射。本面板为机体系投影示意，局部 $+x_b$ 向右，和构型主图的相机左右不等同。$M_z^{(T)}$与$M_y^{(T)}$仅为推力力臂主项；完整映射还含$-\tau_t\sin\delta_t$与$-\tau_f\sin\delta_f$交叉项。零摆角时两台推力均沿 $+x_b$，但前转子 $\bm h_f\parallel +x_b$（CW），尾转子 $\bm h_t\parallel -x_b$（CCW）；两电机反扭矩的轴向投影差形成滚转力矩，虚线为零摆角推力方向。}
\label{fig:thrust_mapping}
\end{figure}

